\section{Shared event mechanisms for the remaining C1 sources}
\label{sec:r6-c1-common-remote}
Normalize the original parent by its positive pure-source coefficient $d$.
Write $D_z=1+z^2$, and write the signed core as $2K$, where
\[
 K=(uW+uvX-Y-vZ)^2+D_u(vY-Z)^2+D_v(X+uZ)^2-D_uD_vZ^2.
 \tag{CE1}
\]
The event packets used below are
\[
4m(vW-X+uvY-uZ)^2,\quad4CD_v(uW-Y)^2,
\quad4AD_uD_vY^2,\quad4\alpha D_u(vY-Z)^2.
\]
Every omitted event and every nonnegative excess is retained as its exact
original square. Source and testing fields remain independent. ``Strong''
means bounded below by a positive real number; ``unbounded'' means larger
than every real number. Neither assertion constrains upper scales.

\begin{theorem}[Two adjacent unbounded events, with complete raw frames]
If $C,A$ are unbounded, the full original parent is tensor SOS at every
projective test for which $u$ is not infinitesimal, and every state.
No physicality assumption beyond nonnegative event weights is needed for
this particular assertion. By the exact cyclic relabeling below, if
$A,\alpha$ are unbounded, the conclusion holds whenever $v$ is not
projectively unbounded.
\end{theorem}
\begin{proof}
Select $C=A=\mu$, with a positive real $\mu$ chosen after the real testing
residues; unboundedness permits every such choice. The following are exact
full quadratic forms, before any restriction to a state ray or contact.
They therefore retain every first state normal and hard-normal coupling.

For $u$ a bounded nonzero unit and $v$ bounded, use the four raw roots
\[
 (W,Y,B,V),\qquad B=X+uZ,\quad V=vZ.
\]
The entire core equals
\[
 K=(uW-Y+uvB)^2+D_vB^2+D_uv^2Y^2+D_uu^2V^2
       -2D_uuV(W+vB).                         \tag{CE2}
\]
The two selected event forms remain $4\mu D_v(uW-Y)^2$ and
$4\mu D_uD_vY^2$. In their common real kernel $W=Y=0$, the core has
matrix on $(B,V)$
\[
 \begin{pmatrix}1+v^2D_u&-D_uuv\\-D_uuv&D_uu^2\end{pmatrix},
 \qquad \det=D_uu^2>0.                       \tag{CE3}
\]
Here and in residue matrices, bounded testing letters mean their real
residues. Thus a sufficiently large finite $\mu$ gives a strictly positive
residue on the entire four-root space, even when $v=0$ literally.

For $p=1/u$ infinitesimal and $v$ bounded, homogenize by $p^2$ and take
\[
 (W,Y,B,R),\qquad B=pX+Z,\quad R=vX.
\]
The exact full core is
\[
 K_p=B^2-2pvBW-2vBY+D_pR^2+2D_pRW
                  +(W-pY)^2+D_pv^2Y^2.       \tag{CE4}
\]
The selected packets are $4\mu D_v(W-pY)^2$ and
$4\mu D_pD_vY^2$. At $p=0$ their sum with $2K_p$ is
\[
 4\mu D_v(W^2+Y^2)+2(W+R)^2+2(B-vY)^2,       \tag{CE5}
\]
strict for every $\mu>0$. The formula allows literal $p=0$ and $v=0$;
the raw roots need not form independent state coordinates.

For projectively unbounded $v$, put $w=1/v$, homogenize by $w^2$, and
use $(W,X,Y,H)$ with $H=wZ$. The exact full core is
\begin{align*}
 K_w={}&u^2H^2-2u^2HY-2uHW+2uwHX+u^2w^2W^2\\
 &+2u^2wWX-2uw^2WY+(D_u+w^2)(X^2+Y^2)-2uwXY.
                                                        \tag{CE6}
\end{align*}
The selected packets are $4\mu D_w(uW-Y)^2$ and
$4\mu D_uD_wY^2$. If $u$ is a bounded nonzero unit, at $w=0$ their
common kernel has $W=Y=0$ and $K_w=D_uX^2+u^2H^2$ strictly positive.
A sufficiently large real $\mu$ again works. If also $p=1/u$ is
infinitesimal, homogenize (CE6) by $p^2$. The exact full core is
\begin{align*}
 K_{pw}={}&H^2-2pHW+2pwHX-2HY+w^2W^2+2wWX-2pw^2WY\\
 &+(D_p+p^2w^2)(X^2+Y^2)-2pwXY.              \tag{CE7}
\end{align*}
At $p=w=0$ the full selected residue is
\[
 4\mu(W^2+Y^2)+2X^2+2(H-Y)^2>0.             \tag{CE8}
\]
All these exact Gram coefficients are bounded. Strictness in (CE3) and
its counterpart following (CE6) implies a finite cap by the elementary
real Schur-complement argument: the event Gram is positive definite on
the complement of its kernel, and the core is positive definite on the
kernel. No source/testing rate comparison is involved. Apply the bounded
positive-residue Gram lemma, and restore every retained original square.

For the second assertion, compose the inherited Alice reflection with
the swap/reflection. The resulting one-step inverse cycle rotation is
\[
 (u',v')=(1/v,u),\quad
 (W',X',Y',Z')=(Y+vZ,-W-vX,Z-vY,vW-X),          \tag{CE9}
\]
and sends event index $i$ to $i-1$ in
$(\lambda,D,C,A,\alpha,m,\beta,n)$. In particular original $A,\alpha$
become new $C,A$. Its full-parent identity is
\[
 v^2F(\mathrm{rot}^{-1}c;1/v,u;\psi')=D_vF(c;u,v;\psi).
                                                        \tag{CE10}
\]
The state map has determinant of absolute value $D_v^2$, and its
homogeneous version at $v=\infty$ is invertible as well. Thus
$u'=1/v$ not infinitesimal is exactly $v$ not projectively unbounded.
Only pure testing factors are restored. Both signs, all literal endpoints,
and complex states (by real/imaginary splitting) are included.
\end{proof}

\begin{theorem}[Four consecutive strong events]
For an original physical parent with $C,A,\alpha,m$ strong, the full
parent is tensor SOS at every test except possibly $u$ a bounded nonzero
unit and $v$ projectively unbounded. All original upper coefficient
scales and all nonnegative unselected coefficients are allowed.
\end{theorem}
\begin{proof}
Cap these four events at any positive real $\mu$ smaller than all four
weights, and retain every excess. Noncommuting-residue tests use the
inherited physical-source theorem, applied to the original full parent.
For commuting-residue tests use the capped parent.

First suppose $v$ is bounded. On the three open commuting axes, the core
and the four caps have zero common kernel: at $u=0$, the core forces
$X=Y=0$, then $m$ and $\alpha$ force $W=Z=0$ when $v$ is a unit;
at $u=\infty$, $C,A$ force $W=Y=0$, then the core forces $X=Z=0$;
at $v=0$, $A,\alpha$ force $Y=Z=0$, then the core forces $W=X=0$
when $u$ is a unit. These are strictly positive bounded state Grams.

At $(u,v)=(0,0)$ the residue is
\[
 (2+4\mu)X^2+(2+8\mu)Y^2+4\mu Z^2.
\]
Only the $W$ ray remains. Normalize $W=1$ and take the five infinitesimal
roots $(u,v,X,Y,Z)$. The full polynomial has no constant or linear term;
its quadratic part is
\[
 2(u-Y)^2+2X^2+
 4\mu\{(v-X)^2+(u-Y)^2+Y^2+Z^2\},            \tag{CE11}
\]
which is strictly positive. At $(u,v)=(\infty,0)$ only the $X$ ray
remains. Normalize $X=1$ and use $(p,v,W,Y,Z)$. Its quadratic part is
\[
 2(W+v)^2+2(p+Z)^2+
 4\mu\{(p+Z)^2+W^2+Y^2+Z^2\},              \tag{CE12}
\]
also strictly positive. All higher monomials in each chart allocate to
products of two displayed roots with infinitesimal bounded coefficients.
This is an exact finite Gram, not a truncation certificate.

For $w=1/v$ infinitesimal use the full raw frame of (CE6). All four
selected packets, with every normal present, are exactly
\[
 4\mu\{(W-wX+uY-uH)^2+D_w(uW-Y)^2
              +D_uD_wY^2+D_u(Y-H)^2\}.        \tag{CE13}
\]
At $u=w=0$ their sum with $2K_w$ has residue
\[
 4\mu W^2+2X^2+(8\mu+2)Y^2+4\mu(Y-H)^2,    \tag{CE14}
\]
strict on all four roots. If $p=1/u$ is infinitesimal, homogenize again;
the exact packets become
\[
 4\mu\{(pW-pwX+Y-H)^2+D_w(W-pY)^2
              +D_pD_wY^2+D_p(Y-H)^2\}.        \tag{CE15}
\]
Together with (CE7), at $p=w=0$ the residue is
\[
 4\mu W^2+2X^2+4\mu Y^2+(8\mu+2)(Y-H)^2,    \tag{CE16}
\]
again strict. These frames require no state normalization, and no
division by either small testing coordinate. Literal $w=0$ is included.
The only chart not covered is ordinary $u$ and infinitesimal $w$.
Restore the original packets and all pure factors as above.
\end{proof}

\begin{lemma}[Matched weak outer events at the large-Bob corner]
Suppose $q>0$ is a pure-source infinitesimal and
\[
 C\ge c q^2,\quad m\ge m_*q^2,\quad
 A\ge a_*/q^2,\quad\alpha\ge b_*/q^2
                                                        \tag{CE17}
\]
for positive real $c,m_*,a_*,b_*$. At projectively unbounded $v$,
the complete parent is tensor SOS whenever $u$ is infinitesimal or
projectively unbounded. In the latter chart the lower bound on $m$
is unnecessary.
\end{lemma}
\begin{proof}
Keep $w=1/v$ and $H=wZ$ as in (CE6)--(CE7). If $u$ is infinitesimal,
use the five raw roots
\[
 E=qW,\quad X,\quad F=Y/q,\quad G=H/q,\quad V=uwW.
\]
The exact whole core (CE6) becomes
\begin{align*}
 K_w={}&V^2+(D_u+w^2)X^2
       +q^2\{u^2(G-F)^2+D_wF^2\}\\
 &-2uGE+2quwX(G-F)+2uVX-2uw^2EF.             \tag{CE18}
\end{align*}
Cap only $m,A,\alpha$ at the indicated weights. Their complete packets
are respectively
\[
 4m_*(E-qwX+uq^2(F-G))^2,\quad
 4a_*D_uD_wF^2,\quad4b_*D_u(F-G)^2.
\]
The selected residue is
\[
 2V^2+2X^2+4m_*E^2+4a_*F^2+4b_*(F-G)^2>0.   \tag{CE19}
\]
Every coefficient in (CE18) is bounded independently of all relative
source/testing rates.

If instead $p=1/u$ is infinitesimal, use
\[
 E=qW,\quad F=Y/q,\quad G=H/q,\quad
 R=X+wW,\quad S=pX.
\]
The exact whole core (CE7) is
\[
 K_{pw}=R^2+D_wS^2+q^2(G-F)^2+p^2D_wq^2F^2
       -2pE(G+w^2F)+2qwS(G-F).               \tag{CE20}
\]
Cap only $C,A,\alpha$. Their complete packets are
\[
 4cD_w(E-pq^2F)^2,\quad4a_*D_pD_wF^2,
 \quad4b_*D_p(F-G)^2.
\]
Together with $2K_{pw}$ their residue is
\[
 2R^2+2S^2+4cE^2+4a_*F^2+4b_*(F-G)^2>0.     \tag{CE21}
\]
Only the pure-source inverse $1/q$ occurs in a root. These exact raw
frames include every state normal and every literal endpoint. The
strict-residue lemma and restoration of excess packets prove the claim.
\end{proof}

\paragraph{Exact application to the original C1 source rows.}
The source premises here are precisely the already proved original
formulas (NC1)--(NC7), (NI3)--(NI10); no independent source family is
substituted. These common theorems give the following new quantified
test reductions, always intersected with prior admissions:
\begin{enumerate}
\item C1.1, every $e$ type: (NI3) gives $C,A$ unbounded;
only $u$ infinitesimal remains.
\item C1.4, $L$ infinitesimal and $\eta$ unbounded: (NI6)
again gives $C,A$ unbounded; only $u$ infinitesimal remains.
\item C1.4, $L$ unbounded: (NI4) and (CE20) leave only
$u$ bounded and $v$ unbounded.
\item C1.5 actual UH, $\eta$ unbounded and $\zeta$ unit:
(NI9) gives four strong events; only $u$ unit and $v$ unbounded remains.
\item C1.5 actual UH, $\zeta$ unbounded: (NI10) and (CE17)
leave the same sole strip, $u$ unit and $v$ unbounded.
\end{enumerate}
For the last line choose $q=\zeta^{-1/2}$. Formula (NI10) gives
$C,m\asymp q^2$, $\alpha\asymp1/(q^2Q)$ and $A\asymp1/r$;
the original $Q=1-x\zeta r>0$, with unit $x$, bounds $Aq^2$ below.
Thus (CE17) holds with positive real constants. The adjacent-event
theorem covers bounded $v$ and (CE18)--(CE21) remove both end types of
$u$ at unbounded $v$. Both actual UH source patterns now have the same
sole untreated testing strip, ordinary $u$ and unbounded $v$.

In the $L$-unbounded C1.4 line take $q=L^{-1/2}$.
Then (NI4) gives $C\asymp q^2$, $\alpha\asymp1/(q^2Q)$ and
$Aq^2\asymp1/(Lr)>1$. The $p$-chart (CE20), which needs no lower
bound on $m$, removes the projectively unbounded-$u$ tests there.

All bounds are relative to original $d$. The four-source-event result
does not require the additional NI9 total-scale condition that at least
one opening be infinitesimal. Every literal permitted $D=0$ or
$\beta=0$ is retained throughout, since neither coefficient was capped.
No claim of nonemptiness or full-row completion is made. In particular,
the $L$ unit and the $L$ infinitesimal/$\eta$ bounded parts of C1.4 are
unchanged, and these remote results do not supply the missing critical
contact conductors.

\paragraph{Scoped corroboration.}
The checker reconstructs the exact core, all six raw identities
(CE2), (CE4), (CE6), (CE7), (CE18), (CE20), their selected packets and residues, the
two corner quadratic forms, and the cyclic core and selected-event
transport. It checks whole quadratic polynomials, hence all state normals
at once. The source entry, ordered-field bounds, real Schur argument,
strict-residue lemma, and the inherited noncommuting theorem remain the
analytical proof. No whole-row count changes follow from these checks.
